\chapter*{编者说明}
\addcontentsline{toc}{chapter}{编者说明}

历史最容易被讲成两种样子：一种是把朝代和皇帝排成名单，另一种是把几个名人的故事讲得热闹。本书都不想做。它想做的事情更朴素：把一个事件放回它发生的那一年，先讲清楚\term{谁做了什么}，再顺着线索往前问“为什么会走到这一步”，往后问“这件事后来改变了什么”。

所以每一章先放一条时间轴，让读者知道自己站在什么位置；随后用短条目推进；真正改变局势的大事，再停下来讲它前面的积累、当时的人怎么选、以及它给后世留下了什么。相邻的两件事不必然互为因果，后人看到的结果也未必是当事人最初想要的结果，本书会把这些区别说出来。

叙述从周开始。西周早期的绝对年份并不总是没有争议，书中会直接标出来；自公元前841年共和元年起，传统编年有了连续而稳定的锚点，条目会随之变得更密。底层史料以《史记》、二十四史和《资治通鉴》为主线，同时提醒读者：它们不是透明的录像，而是不同时代的人留下的叙事与判断。
